\chapter{Wstęp}

% Miejsce na podrozdziały 1.1 i 1.2 (np. Wprowadzenie)

\section{Przegląd aktualnych rozwiązań}

Obecnie najpowszechniejszym rozwiązaniem wykorzystywanym na stanowiskach pracy są klasyczne lampy biurkowe, które opierają się na układach wymagających ingerencji użytkownika. Zaletą tradycyjnych urządzeń oświetleniowych jest niska złożoność i prostota użycia.

Z kolei do głównych minusów obecnych, klasycznych rozwiązań należą:
\begin{itemize}
\item Brak automatycznego dostosowania oświetlenia wewnętrznego do natężenia światła słonecznego lub zewnętrznego \cite{karta_projektu}.
\item Ręczna regulacja jasności, co bywa wysoce niepraktyczne, zwłaszcza w środowisku i miejscach pracy, gdzie warunki oświetleniowe podlegają częstym zmianom \cite{karta_projektu}.
\item Powodowanie nadmiernego zmęczenia wzroku u osób pracujących przy komputerze \cite{karta_projektu}.
\item Niepotrzebne zużycie energii elektrycznej w momentach, gdy światło z zewnątrz jest wystarczające \cite{karta_projektu}.
\end{itemize}